% Usage in text:
% \ac{<label>}   Full form (long (short)) on first use, short form afterwards.
% \acs{<label>}  Short form only.
% \acf{<label>}  Always the full form.
% \acl{<label>}  Long form only.
% Plural: append "p", e.g. \acp{VLM}. Capitalised: \Ac, \Acs, \Acf, \Acl.
%
% Long forms are stored in lower case to ensure that \ac{} reads correctly inside a
% sentence; use \Ac{} at the beginning of a sentence. Only abbreviations that
% are actually used in the text appear in the List of Abbreviations.
%
% Use {} brackets if a value contains special characters or a comma!

\newabbreviation{AoV}{AoV}{angle of view}

\newabbreviation{BAG}{BAG}{Basisregistratie Adressen en Gebouwen}

\newabbreviation{CNN}{CNN}{convolutional neural network}

\newabbreviation{CV}{CV}{computer vision}

\newabbreviation{EPC}{EPC}{Energy Performance Certificate}

\newabbreviation{LightGBM}{LightGBM}{Light Gradient Boosting Machine}

\newabbreviation{LoD}{LoD}{Level of Detail}

\newabbreviation{ML}{ML}{machine learning}

\newabbreviation{MLP}{MLP}{multilayer perceptron}

\newabbreviation{OSM}{OSM}{OpenStreetMap}

\newabbreviation{SVI}{SVI}{street view imagery}

\newabbreviation{TABULA}{TABULA}{Typology Approach for Building Stock Energy Assessment}

\newabbreviation{UBEM}{UBEM}{urban building energy modelling}

\newabbreviation{ViT}{ViT}{Vision Transformer}

\newabbreviation{VLM}{VLM}{vision-language model}
